%%
%% ACM Multimedia 2026 — Supplementary Material
%%
\documentclass[sigconf,review,anonymous]{acmart}

\usepackage{amsmath,amssymb,amsfonts}
\usepackage{booktabs}
\usepackage{multirow}
\usepackage{graphicx}
\usepackage{xcolor}
\usepackage{bm}
\usepackage{enumitem}
\usepackage{pifont}

\copyrightyear{2026}
\acmYear{2026}
\setcopyright{none}

\newcommand{\diagonal}{\textsc{Span}}
\newcommand{\stem}{\textsc{Stem}}
\newcommand{\sdiag}{\mathcal{S}_{\mathrm{diag}}}
\newcommand{\Mobs}{M_{\mathrm{obs}}}
\newcommand{\Mideal}{M_{\mathrm{ideal}}}
\newcommand{\calB}{\mathcal{B}}
\newcommand{\calT}{\mathcal{T}}
\newcommand{\calP}{\mathcal{P}}
\newcommand{\calH}{\mathcal{H}}
\newcommand{\Eij}{E_{ij}}

\newtheorem{lemma}{Lemma}

\begin{document}

\title{Supplementary Material for\\The Algebra of Storytelling: A State Transition Benchmark\\for Spanning Multi-Entity Multi-Shot Narratives}

\author{Anonymous Authors}
\affiliation{\institution{Anonymous Institution}}

\maketitle

\appendix
\renewcommand{\thetheorem}{\arabic{theorem}}
\renewcommand{\thelemma}{\arabic{lemma}}

%% ============================================================
\section{Necessity of the Split Basis ($\calB_2$)}
\label{app:necessity_b2}
%% ============================================================

In Theorem~1 (main text), we established that $\{\calB_1, \calB_2\}$ spans $\mathbb{M}_{2 \times 3}(\mathbb{Z})$. To avoid redundancy with the main text's constructive proof, we provide the formal algebraic proof of why $\calB_1$ (Relay) alone is strictly insufficient to span the space, thereby necessitating the introduction of $\calB_2$ (Split).

\begin{lemma}[Row-Sum Invariance of $\calB_1$]
Let $R_i(M)$ denote the sum of elements in the $i$-th row of a matrix $M$. For any character role swapping $\sigma \in \Sigma_3$, the permuted Relay matrix $M = \sigma \circ \calB_1$ strictly preserves the topological invariant:
\begin{equation}
R_1(M) = +1, \quad R_2(M) = -1.
\end{equation}
\end{lemma}

\begin{proof}
The canonical Relay matrix $\calB_1 = \bigl[\begin{smallmatrix} 0 & 1 & 0 \\ -1 & 0 & 0 \end{smallmatrix}\bigr]$ trivially satisfies $R_1(\calB_1) = 1$ and $R_2(\calB_1) = -1$. Because the operator $\sigma$ solely permutes column indices, it strictly preserves the algebraic sum of each row.
\end{proof}

Consequently, any $\mathbb{Z}$-linear combination derived exclusively from Relay variants, $X = \sum_k \alpha_k (\sigma_k \circ \calB_1)$, will possess row sums $R_1(X) = \sum_k \alpha_k$ and $R_2(X) = -\sum_k \alpha_k$, imposing a strict global constraint:
\begin{equation}
R_1(X) + R_2(X) = 0.
\end{equation}
However, the standard elementary matrix $E_{11} = \bigl[\begin{smallmatrix} 1 & 0 & 0 \\ 0 & 0 & 0 \end{smallmatrix}\bigr]$ yields $R_1(E_{11}) + R_2(E_{11}) = 1 \neq 0$. Thus, $\calB_1$ variants alone cannot break this symmetric invariance. The Split matrix $\calB_2 = \bigl[\begin{smallmatrix} 0 & -1 & 0 \\ -1 & 1 & 0 \end{smallmatrix}\bigr]$, which yields $R_1(\calB_2) + R_2(\calB_2) = -1$, is therefore algebraically essential to span the full narrative space.


%% ============================================================
%% ============================================================
%% ============================================================
\section{Generalization to Arbitrary Narrative Scales}
%% ============================================================
\label{app:generalization}
%% ============================================================

In Theorem~1 (main text), we proved completeness for the canonical 3-shot, 3-entity subspace ($\mathbb{M}_{2\times 3}(\mathbb{Z})$). We now construct a formal algebraic pipeline to generalize this spanning property to arbitrary narrative scales: expanding first spatially to $N$ entities, then temporally to $T$ shots.
%% ============================================================
\section{Generalization to Arbitrary Narrative Scales}
\label{app:generalization}
%% ============================================================

In Theorem~1 (main text), we proved completeness for the canonical 3-shot, 3-entity subspace ($\mathbb{M}_{2\times 3}(\mathbb{Z})$). We now construct a formal algebraic pipeline to generalize this spanning property to arbitrary narrative scales: expanding first spatially to $N$ entities, then temporally to $T$ shots.

\subsection{Phase 1: Spatial Expansion (Arbitrary Entities, $N \geq 3$)}

To accommodate an arbitrary cast size, we define the \textbf{Spatial Zero-Padding Operator ($\mathcal{Z}_N$)}, which appends $N-3$ zero-columns to a canonical local matrix:
\begin{flalign}
& \mathcal{Z}_N(X) = \bigl[\; X \;\;\; \mathbf{0}_{2\times(N-3)} \;\bigr], \quad X \in \mathbb{M}_{2\times 3}. &&
\end{flalign}

% ==========
Similar to temporal embedding, $\mathcal{Z}_N$ operates as a strict linear homomorphism:
\begin{equation}
\mathcal{Z}_N(c_1 X + c_2 Y) = c_1\mathcal{Z}_N(X) + c_2\mathcal{Z}_N(Y), \quad \forall c_1, c_2 \in \mathbb{Z}.
\end{equation}
This linearity ensures that any algebraic superposition constructed within the canonical local space is perfectly preserved upon spatial expansion.
% ======================================

We define the spatially padded bases as $\hat{\calB}_m = \mathcal{Z}_N(\calB_m)$ for $m \in \{1, 2\}$. 

Because any target entity $e \in \{1, \dots, N\}$ can be routed to the active first three columns via the extended permutation group $\Sigma_N$, every elementary matrix $E^{(\text{local})}_{i,e} \in \mathbb{M}_{2 \times N}(\mathbb{Z})$ is algebraically constructible from the canonical base $E_{i,j} \in \mathbb{M}_{2 \times 3}$:
\begin{flalign}
& E^{(\text{local})}_{i,e} = \sigma \circ \mathcal{Z}_N \bigl( E_{i,j} \bigr), \quad \exists \sigma \in \Sigma_N, \; j \in \{1,2,3\}. &&
\end{flalign}
Consequently, the spatially expanded bases seamlessly span the entire multi-entity local space:
\begin{flalign}
& \mathbb{M}_{2 \times N}(\mathbb{Z}) = \operatorname{span}_{\mathbb{Z}}\bigl(\{ \sigma \circ \hat{\calB}_m \mid m \in \{1,2\}, \sigma \in \Sigma_N \}\bigr). &&
\label{eq:spatial_span}
\end{flalign}

\subsection{Phase 2: Temporal Expansion (Arbitrary Shots, $T \geq 3$)}

To map our local multi-entity patterns into a continuous $T$-shot timeline, we define the \textbf{Temporal Shift Operator ($\mathcal{S}_\tau$)}. This operator embeds a 2-step transition block into the global temporal window starting at $\tau \in \{1, \dots, T-2\}$:
\begin{flalign}
& \mathcal{S}_\tau(X) = \begin{bmatrix} \mathbf{0}_{(\tau-1)\times N} \\ X \\ \mathbf{0}_{(T-\tau-2)\times N} \end{bmatrix}, \quad X \in \mathbb{M}_{2\times N}. &&
\end{flalign}
Crucially, $\mathcal{S}_\tau$ strictly satisfies linear homomorphism:
\begin{equation}
\mathcal{S}_\tau(c_1 X + c_2 Y) = c_1\mathcal{S}_\tau(X) + c_2\mathcal{S}_\tau(Y), \quad \forall c_1, c_2 \in \mathbb{Z}.
\end{equation}
This linearity structurally guarantees that the superposition of local transitions maps directly to the global space without cross-temporal interference.

\subsection{Phase 3: The Grand Spanning Theorem}

By synthesizing the spatial and temporal expansions, we formalize the ultimate completeness of the \diagonal{} framework.

\setcounter{theorem}{1}
\begin{theorem}[Grand Topological Generalization]
For any narrative of arbitrary length ($T \geq 3$) involving an arbitrary cast size ($N \geq 3$), every conceivable state transition matrix $\Delta S \in \mathbb{M}_{(T-1)\times N}(\mathbb{Z})$ can be algebraically generated by:
\begin{equation}
\Delta S = \sum_{\tau=1}^{T-2} \sum_{m=1}^{2} c_{\tau,m}\; \bigl(\mathcal{S}_\tau \circ \sigma_{\tau,m} \circ \mathcal{Z}_N\bigr)(\calB_m),
\label{eq:grand_spanning}
\end{equation}
where $c_{\tau,m} \in \mathbb{Z}$ and $\sigma_{\tau,m} \in \Sigma_N$.
\end{theorem}

\begin{proof}
Let $E^{(\text{global})}_{t,e} \in \mathbb{M}_{(T-1)\times N}(\mathbb{Z})$ denote the standard global elementary matrix. Since any $\Delta S$ is a $\mathbb{Z}$-linear combination of $\{E^{(\text{global})}_{t,e}\}$, it suffices to show that every such elementary matrix is constructible.

For any target transition step $t \in \{1, \dots, T-1\}$, we define an exact index mapping to isolate the active local window $\tau$ and the local row index $i \in \{1, 2\}$:
\begin{equation}
\tau = \min(t, \, T-2), \quad i = t - \tau + 1.
\label{eq:index_mapping}
\end{equation}
By Eq.~\ref{eq:spatial_span}, the spatial operations can generate the required local matrix $E^{(\text{local})}_{i,e}$. Applying the linear shift operator $\mathcal{S}_\tau$ precisely projects this local block to the target global row $t$:
\begin{equation}
E^{(\text{global})}_{t,e} = \mathcal{S}_\tau \bigl( E^{(\text{local})}_{i,e} \bigr) \;\in\; \operatorname{span}_{\mathbb{Z}}\bigl( \{ \mathcal{S}_\tau \circ \sigma \circ \hat{\calB}_m \} \bigr).
\end{equation}
Thus, all standard bases in the global domain are generated, mathematically certifying complete global coverage. \qedhere
\end{proof}

\paragraph{Dimensionality and Scalability.}
The target global narrative space has a formidable dimensionality:
\begin{equation}
\dim\bigl(\mathbb{M}_{(T-1) \times N}(\mathbb{Z})\bigr) = (T-1) \cdot N.
\end{equation}
However, Eq.~\ref{eq:grand_spanning} proves that this space is structurally spanned by sliding just $2$ base matrices ($\calB_1, \calB_2$) across $(T-2)$ temporal windows. This mathematical guarantee ensures that narrative evaluation scales \emph{linearly} as $\mathcal{O}(T)$, completely bypassing the intractable combinatorial explosion $\mathcal{O}(2^{TN})$ associated with empirical prompt-sampling methodologies.
%% ============================================================
\section{Dataset Construction Details}
\label{app:dataset}
%% ============================================================

\begin{table}[ht]
\caption{Complete benchmark composition (1,000 scenarios, 3,000 shots).}
\label{tab:full_composition}
\centering
\small
\begin{tabular}{@{}llccc@{}}
\toprule
Type & Pattern & Formula & Scenarios & Shots \\
\midrule
\multirow{3}{*}{Basis}
  & Relay         & $A \to AB \to B$     & 200 & 600 \\
  & Split         & $AB \to A \to B$     & 200 & 600 \\
  & Accumulation  & $A \to AB \to ABC$   & 200 & 600 \\
\midrule
\multirow{4}{*}{Derived}
  & Convergence      & $A \to B \to AB$     & 100 & 300 \\
  & Sliding Window   & $AB \to BC \to C$    & 100 & 300 \\
  & Reduction        & $ABC \to AB \to A$   & 100 & 300 \\
  & Reverse Relay    & $B \to AB \to A$     & 100 & 300 \\
\midrule
\multicolumn{2}{l}{Total} & & 1,000 & 3,000 \\
\bottomrule
\end{tabular}
\end{table}

Table~\ref{tab:full_composition} summarizes the benchmark composition.

\paragraph{Prompt generation.}
All prompts are generated via Google Gemini~1.5 (temperature 0.7) using structured templates specifying narrative pattern, entity descriptions, and scene backgrounds.
Of 1,200 initial candidates, 200 (16.7\%) were discarded for ambiguous descriptions, visually indistinguishable entity pairs, or degenerate backgrounds.

\paragraph{Scene themes.}
10 environmental themes: Domestic\_LivingRoom, Domestic\_Kitchen, Urban\_Caf\'{e}, Urban\_Library, Urban\_Street, Nature\_Forest, Nature\_Beach, Nature\_Park, Professional\_Office, Professional\_Hospital.

\paragraph{Entity type distribution.}
Three categories: Human (H), Professional (P), Animal (A). Table~\ref{tab:entity_types} shows the full distribution.

\begin{table}[ht]
\caption{Entity type combination distribution.}
\label{tab:entity_types}
\centering
\small
\begin{tabular}{@{}lcc@{}}
\toprule
Entity Types & Count & \% \\
\midrule
(H, P)    & 268 & 26.8 \\
(H, P, P) & 169 & 16.9 \\
(P, P)    & 110 & 11.0 \\
(H, H)    & 96  & 9.6 \\
(H, H, P) & 93  & 9.3 \\
(A, H)    & 72  & 7.2 \\
(A, A)    & 54  & 5.4 \\
(P, P, P) & 45  & 4.5 \\
(A, A, A) & 38  & 3.8 \\
(A, A, H) & 35  & 3.5 \\
(H, H, H) & 20  & 2.0 \\
\bottomrule
\end{tabular}
\end{table}


%% ============================================================
\section{Prompt Quality Control}
\label{app:quality_control}
%% ============================================================

All prompts undergo automated validation ensuring entity--pattern topology consistency: the set of entities mentioned in each shot description must exactly match the presence matrix $S$ for the target pattern.
Additionally, 10\% of prompts are manually reviewed by two annotators (97\% adherence to intended structure, Cohen's $\kappa = 0.91$).
Prompts that describe ambiguous scenarios (e.g., entities that could plausibly be interpreted as present in a scene's background) are revised for clarity.


%% ============================================================
\section{Sample Prompts}
\label{app:prompts}
%% ============================================================

Each prompt has two variants: \texttt{gen\_prompt} (+ quality tags for generation) and \texttt{eval\_prompt} (clean scene descriptions for evaluation). Entity names are extracted directly from the eval prompts for VLM entity queries. Tables~\ref{tab:sample_relay}--\ref{tab:sample_slidwin} show representative examples for three patterns.

\begin{table}[ht]
\caption{Relay pattern sample ($A \to AB \to B$).}
\label{tab:sample_relay}
\centering
\small
\setlength{\tabcolsep}{3pt}
\begin{tabular}{@{}clp{5.5cm}@{}}
\toprule
Shot & Field & Prompt \\
\midrule
\multirow{2}{*}{1}
  & eval & A curious toddler stacking colorful wooden blocks on a soft rug in a sunlit living room \\
  & entity &a toddler \\
\midrule
\multirow{2}{*}{2}
  & eval & A toddler and a mother sitting together on the rug, the mother helping stack blocks \\
  & entity &a toddler and a mother together \\
\midrule
\multirow{2}{*}{3}
  & eval & A mother tidying up scattered blocks on the rug, the toddler no longer present \\
  & entity &a mother \\
\bottomrule
\end{tabular}
\end{table}

\begin{table}[ht]
\caption{Accumulation sample ($A \to AB \to ABC$).}
\label{tab:sample_accum}
\centering
\small
\setlength{\tabcolsep}{3pt}
\begin{tabular}{@{}clp{5.5cm}@{}}
\toprule
Shot & Field & Prompt \\
\midrule
1 & eval & A lone eagle soaring above a dense forest canopy at dawn \\
  & entity &an eagle \\
\midrule
2 & eval & An eagle and a fox standing at the edge of a forest clearing \\
  & entity &an eagle and a fox \\
\midrule
3 & eval & An eagle perched on a branch, a fox below, and a rabbit emerging from undergrowth \\
  & entity &an eagle, a fox, and a rabbit \\
\bottomrule
\end{tabular}
\end{table}

\begin{table}[ht]
\caption{Sliding Window sample ($AB \to BC \to C$).}
\label{tab:sample_slidwin}
\centering
\small
\setlength{\tabcolsep}{3pt}
\begin{tabular}{@{}clp{5.5cm}@{}}
\toprule
Shot & Field & Prompt \\
\midrule
1 & eval & A chef and a waiter preparing a large banquet table in an elegant restaurant \\
  & entity &a chef and a waiter \\
\midrule
2 & eval & A waiter greeting a sommelier at the restaurant entrance, the chef no longer visible \\
  & entity &a waiter and a sommelier \\
\midrule
3 & eval & A sommelier examining a wine bottle under warm chandelier light, alone \\
  & entity &a sommelier \\
\bottomrule
\end{tabular}
\end{table}


%% ============================================================
\section{Model Implementation Details}
\label{app:models}
%% ============================================================

\begin{table}[ht]
\caption{Inference specifications. $^\dagger$T2I: 30 steps; I2V: 50 steps.}
\label{tab:inference}
\centering
\small
\setlength{\tabcolsep}{3pt}
\begin{tabular}{@{}lllllc@{}}
\toprule
Model & Backbone & Paradigm & Steps & Resolution & Frames \\
\midrule
StoryDiff & SDXL+SEINE & T2I$\to$I2V$^\dagger$ & 30+50 & 512$\times$512 & 16 \\
EchoShot  & Wan2.1-1.3B & T2V (FM) & 50 & 832$\times$480 & 93 \\
VIC       & CogVideoX-5B & In-context & 50 & 720$\times$480 & 49 \\
VGoT      & Kolors+DC & T2I$\to$I2V & 25+50 & 512$\times$512 & 16 \\
\bottomrule
\end{tabular}
\end{table}

Table~\ref{tab:inference} summarizes the inference configurations. All models use a fixed seed, negative prompt: ``low quality, blurry, deformed, morphed, distorted, extra limbs, bad anatomy, text, watermark''.
All four models successfully generated 1000 videos each (4,000 total).


%% ============================================================
\section{Evaluation Details}
\label{app:eval_details}
%% ============================================================

\paragraph{VLM entity detector.}
Qwen2-VL-7B-Instruct~\cite{qwen2vl} (\texttt{Qwen/Qwen2-VL-7B-Instruct}), fp16 precision, greedy decoding (\texttt{do\_sample=False, max\_new\_tokens=5}).
One middle frame per shot is extracted ($T = 3$ shots per video, 3 frames total).

\paragraph{Entity query.}
Binary yes/no: ``Is there a \{entity\} visible in this image? Answer only Yes or No.'' The VLM response is parsed by prefix: if it starts with ``yes'' (case-insensitive), $S_\text{obs}[t,k] = 1$; otherwise 0. No thresholds are required.

\paragraph{Detection quality.}
Across 4,000 videos (28,800 cells): TPR $= 52.1\%$, TNR $= 66.7\%$.
The low TPR reflects entity ambiguity in generated images rather than VLM failure---manual inspection of 100 randomly sampled false-negative cells confirmed that 65\% were genuine generation failures, 19\% were ambiguous, and only 16\% were VLM detection errors, yielding 80.2\% VLM precision (see main paper, \S5.6). The same VLM achieves 100\% accuracy on unambiguous GT cat videos (Appendix~\ref{app:gt_failures}).

\paragraph{Failure classification.}
Given $E_{t,e} = \Delta S_{\text{obs},t,e} - \Delta S^*_{t,e}$:
\begin{itemize}[leftmargin=*,nosep]
  \item Type~I (Context Bleeding): $\Delta S^*_{t,e} = -1$ and $E_{t,e} = +1$ (entity should exit but persists).
  \item Type~II (Amnesia): $\Delta S^*_{t,e} = +1$ and $E_{t,e} = -1$ (entity should enter but absent).
  \item Type~III (Spurious): $\Delta S^*_{t,e} = 0$ and $E_{t,e} \neq 0$ (unprescribed transition).
  \item Type~IV (Reversal): $|E_{t,e}| = 2$ (opposite direction; strict subset of I/II).
\end{itemize}


%% ============================================================
\section{State Transition Matrices Library}
\label{app:delta_s}
%% ============================================================

Table~\ref{tab:delta_s} lists the prescribed transition matrices $\Delta S^*$ and presence matrices $S^*$ for all 7 canonical patterns.

\begin{table}[ht]
\caption{Prescribed transition matrices ($\Delta S^*$) and presence matrices ($S^*$) for all 7 canonical patterns in MEMS-1K.}
\label{tab:delta_s}
\centering
\small
\renewcommand{\arraystretch}{1.2}
\setlength{\tabcolsep}{8pt}
\begin{tabular}{@{}lcc@{}}
\toprule
Pattern & Transition $\Delta S^*$ & Presence $S^*$ \\
\midrule
Relay &
$\begin{bmatrix} 0 & 1 \\ -1 & 0 \end{bmatrix}$ &
$\begin{bmatrix} 1 & 0 \\ 1 & 1 \\ 0 & 1 \end{bmatrix}$ \\[12pt]

Split &
$\begin{bmatrix} 0 & -1 \\ -1 & 1 \end{bmatrix}$ &
$\begin{bmatrix} 1 & 1 \\ 1 & 0 \\ 0 & 1 \end{bmatrix}$ \\[12pt]

Accumulation &
$\begin{bmatrix} 0 & 1 & 0 \\ 0 & 0 & 1 \end{bmatrix}$ &
$\begin{bmatrix} 1 & 0 & 0 \\ 1 & 1 & 0 \\ 1 & 1 & 1 \end{bmatrix}$ \\[12pt]

Convergence &
$\begin{bmatrix} -1 & 1 \\ 1 & 0 \end{bmatrix}$ &
$\begin{bmatrix} 1 & 0 \\ 0 & 1 \\ 1 & 1 \end{bmatrix}$ \\[12pt]

Sliding Window &
$\begin{bmatrix} -1 & 0 & 1 \\ 0 & -1 & 0 \end{bmatrix}$ &
$\begin{bmatrix} 1 & 1 & 0 \\ 0 & 1 & 1 \\ 0 & 0 & 1 \end{bmatrix}$ \\[12pt]

Reduction &
$\begin{bmatrix} 0 & 0 & -1 \\ 0 & -1 & 0 \end{bmatrix}$ &
$\begin{bmatrix} 1 & 1 & 1 \\ 1 & 1 & 0 \\ 1 & 0 & 0 \end{bmatrix}$ \\[12pt]

Reverse Relay &
$\begin{bmatrix} 1 & 0 \\ 0 & -1 \end{bmatrix}$ &
$\begin{bmatrix} 0 & 1 \\ 1 & 1 \\ 1 & 0 \end{bmatrix}$ \\
\bottomrule
\end{tabular}
\end{table}


%% ============================================================
\section{Full Per-Pattern Results}
\label{app:per_pattern}
%% ============================================================

\begin{table*}[ht]
\caption{Complete per-pattern results (VLM-based). Best per pattern in \textbf{bold}.}
\label{tab:full_pattern}
\centering
\small
\setlength{\tabcolsep}{3pt}
\begin{tabular}{@{}l|cc|cc|cc|cc@{}}
\toprule
& \multicolumn{2}{c|}{StoryDiffusion} & \multicolumn{2}{c|}{EchoShot} & \multicolumn{2}{c|}{VGoT} & \multicolumn{2}{c}{VIC} \\
Pattern & PTM & EPA & PTM & EPA & PTM & EPA & PTM & EPA \\
\midrule
Relay (B, 2-ent)       & \textbf{.190} & .536 & .138 & .506 & .085 & \textbf{.531} & .080 & .522 \\
Split (B, 2-ent)       & \textbf{.233} & \textbf{.553} & .165 & .538 & .108 & .537 & .048 & .497 \\
Accum.\ (B, 3-ent)    & \textbf{.253} & \textbf{.662} & .095 & .610 & .070 & .642 & .072 & .594 \\
Converg.\ (D, 2-ent)  & \textbf{.277} & \textbf{.595} & .187 & .558 & .053 & .497 & .067 & .555 \\
Slid.\ Win.\ (D, 3-ent) & \textbf{.357} & \textbf{.627} & .270 & .611 & .103 & .502 & .077 & .523 \\
Reduction (D, 3-ent)   & \textbf{.250} & .644 & .195 & .601 & .110 & .588 & .080 & \textbf{.618} \\
Rev.\ Relay (D, 2-ent) & .200 & \textbf{.562} & \textbf{.225} & .555 & .150 & .537 & .080 & .540 \\
\bottomrule
\end{tabular}
\end{table*}

Table~\ref{tab:full_pattern} reports the complete per-pattern results. Sliding Window produces the widest PTM spread (0.077--0.357), making it the most discriminative test.
VIC's PTM range is only 0.032 (0.048--0.080), confirming a fixed coherence ceiling invariant to pattern complexity.


%% ============================================================
\section{Human Evaluation Protocol}
\label{app:human_eval}
%% ============================================================

\paragraph{Platform.}
Custom Streamlit web application showing two anonymized model outputs side by side (3 keyframes each). Entity presence/absence annotations shown below. Left/right randomized.

\paragraph{Scale.}
5-point Likert: 1 = ``A clearly better'', 3 = ``Similar'', 5 = ``B clearly better''.
Two dimensions: Narrative Compliance and Visual Quality, assessed independently.

\paragraph{Criteria.}
\emph{Narrative}: ``Which video better follows the scripted entity dynamics?''
\emph{Visual}: ``Which video looks more visually coherent and natural? Ignore narrative correctness.''

\paragraph{Design.}
120 pairwise comparisons (40 scenarios $\times$ 6 model pairs). Each evaluator assessed 30 pairs via rotating-window scheme (3 raters per pair). 15 graduate student participants (CS and media studies). Median time: $\sim$25 min. Two attention-check trials (100\% pass rate). Table~\ref{tab:score_dist} shows score distributions, and Table~\ref{tab:head_to_head} reports head-to-head preferences.

\begin{table}[ht]
\caption{Score distributions by dimension.}
\label{tab:score_dist}
\centering
\small
\begin{tabular}{@{}lccccccc@{}}
\toprule
& \multicolumn{5}{c}{Score \%} & \multicolumn{2}{c}{Summary} \\
Dim. & 1 & 2 & 3 & 4 & 5 & $\mu$ & $\sigma$ \\
\midrule
Narr.  & 12.3 & 14.8 & 42.3 & 17.1 & 13.5 & 3.05 & 1.18 \\
Visual & 16.7 & 18.2 & 26.0 & 22.4 & 16.7 & 3.04 & 1.34 \\
\bottomrule
\end{tabular}
\end{table}

\begin{table}[ht]
\caption{Head-to-head preferences (\%). M1 win / Tie / M2 win.}
\label{tab:head_to_head}
\centering
\small
\setlength{\tabcolsep}{3pt}
\begin{tabular}{@{}lcc@{}}
\toprule
Matchup & Narrative & Visual \\
\midrule
SD vs.\ Echo  & 32 / 43 / 25 & 28 / 31 / 41 \\
SD vs.\ VIC   & 38 / 24 / 38 & 21 / 33 / 46 \\
SD vs.\ VGoT  & 35 / 40 / 25 & 30 / 28 / 42 \\
Echo vs.\ VIC & 42 / 38 / 20 & 22 / 27 / 51 \\
Echo vs.\ VGoT & 37 / 42 / 21 & 31 / 30 / 39 \\
VIC vs.\ VGoT & 25 / 45 / 30 & 43 / 25 / 32 \\
\bottomrule
\end{tabular}
\end{table}


%% ============================================================
\section{VLM Detection Quality by Pattern}
\label{app:vlm_quality}
%% ============================================================

The VLM detector's TPR and TNR vary across patterns, reflecting differences in entity count and scene complexity (Table~\ref{tab:vlm_quality}).

\begin{table}[ht]
\caption{VLM detection quality (TPR/TNR) by narrative pattern.}
\label{tab:vlm_quality}
\centering
\small
\begin{tabular}{@{}lccc@{}}
\toprule
Pattern & Entities & TPR & TNR \\
\midrule
Relay         & 2 & 51.7\% & 53.6\% \\
Split         & 2 & 51.6\% & 56.2\% \\
Accumulation  & 3 & 52.6\% & 82.9\% \\
Convergence   & 2 & 54.9\% & 55.5\% \\
Sliding Win.  & 3 & 45.6\% & 70.3\% \\
Reduction     & 3 & 54.2\% & 75.4\% \\
Rev.\ Relay   & 2 & 54.1\% & 56.2\% \\
\bottomrule
\end{tabular}
\end{table}

3-entity patterns (Accumulation, Sliding Window, Reduction) achieve substantially higher TNR (70--83\%) because more entities are absent per shot, and the VLM correctly identifies absent entities.
TPR is relatively stable across patterns (46--55\%), reflecting the fundamental difficulty of detecting entities in generated images.
The VLM-based binary detection avoids the threshold tuning required by continuous CLIP probes, making PTM and EPA threshold-free metrics.


%% ============================================================
\section{Entry vs.\ Exit Operation Accuracy}
\label{app:entry_exit}
%% ============================================================

We decompose PTM into entry (prescribed $\Delta S^* = +1$) and exit (prescribed $\Delta S^* = -1$) operations (Table~\ref{tab:entry_exit}).

\begin{table}[ht]
\caption{PTM decomposed by operation type. $\Delta$ = entry $-$ exit accuracy.}
\label{tab:entry_exit}
\centering
\small
\begin{tabular}{@{}lcccc@{}}
\toprule
Model & Entry (+1) & Exit ($-1$) & $\Delta$ & $n$ \\
\midrule
SD    & \textbf{.256} & \textbf{.242} & +.014 & 2,400 \\
Echo  & .154 & .190 & $-$.036 & 2,400 \\
VGoT  & .068 & .120 & $-$.052 & 2,400 \\
VIC   & .068 & .069 & $-$.001 & 2,400 \\
\bottomrule
\end{tabular}
\end{table}

\textbf{VGoT has an entry-exit asymmetry} ($\Delta = -0.052$): exit accuracy (0.120) is nearly double entry accuracy (0.068), confirming the planning-rendering bottleneck where DynamiCrafter cannot introduce new entities absent from the keyframe.
StoryDiffusion is the only model where entry accuracy slightly exceeds exit accuracy ($\Delta = +0.014$), suggesting CSA's shared attention can introduce new entities.
VIC is approximately symmetric ($\Delta \approx 0$) but both entry and exit accuracy are near-chance (0.068--0.069).


%% ============================================================
\section{2-Entity vs.\ 3-Entity Stratified Analysis}
\label{app:entity_count}
%% ============================================================

Patterns use either 2 entities (Relay, Split, Convergence, Reverse Relay; $n = 600$ per model) or 3 entities (Accumulation, Sliding Window, Reduction; $n = 400$ per model). Table~\ref{tab:entity_strat} reports the stratified results.

\begin{table}[ht]
\caption{Metrics stratified by entity count.}
\label{tab:entity_strat}
\centering
\small
\setlength{\tabcolsep}{3pt}
\begin{tabular}{@{}l cccc@{}}
\toprule
& \multicolumn{2}{c}{2-Entity} & \multicolumn{2}{c}{3-Entity} \\
\cmidrule(lr){2-3} \cmidrule(lr){4-5}
Model & PTM & EPA & PTM & EPA \\
\midrule
SD    & .221 & .556 & \textbf{.278} & \textbf{.649} \\
Echo  & .169 & .534 & .164 & .608 \\
VGoT  & .098 & .528 & .088 & .593 \\
VIC   & .067 & .522 & .075 & .582 \\
\bottomrule
\end{tabular}
\end{table}

3-entity patterns achieve higher EPA across all models (0.582--0.649 vs.\ 0.522--0.556) due to the larger denominator reducing per-cell noise impact.
PTM rankings are preserved in both strata: SD $>$ Echo $>$ VGoT $>$ VIC for 2-entity, and SD $>$ Echo $>$ VGoT $\approx$ VIC for 3-entity.
StoryDiffusion benefits most from 3-entity patterns (PTM: 0.221 $\to$ 0.278), suggesting CSA excels when more entities provide stronger discriminative anchors.


%% ============================================================
\section{EPA-PTM Correlation Analysis}
\label{app:correlation}
%% ============================================================

\begin{table}[ht]
\caption{EPA-PTM correlation. Per-video Pearson $r$.}
\label{tab:correlation}
\centering
\small
\begin{tabular}{@{}lcc@{}}
\toprule
Scope & Pearson $r$ & $p$ \\
\midrule
SD only     & .661 & $1.2 \times 10^{-126}$ \\
Echo only   & .530 & $1.9 \times 10^{-73}$ \\
VGoT only   & .472 & $1.1 \times 10^{-56}$ \\
VIC only    & .452 & $1.6 \times 10^{-51}$ \\
All pooled  & .542 & $<10^{-300}$ \\
\bottomrule
\end{tabular}
\end{table}

As shown in Table~\ref{tab:correlation}, EPA and PTM show strong positive correlation at the per-video level ($r = 0.542$ pooled), confirming they measure related but distinct aspects of narrative fidelity.
The correlation is strongest for StoryDiffusion ($r = 0.661$) because it has the widest PTM range, providing more variance for correlation.
All per-method correlations are highly significant ($p < 10^{-51}$), in contrast to the CLIP-based analysis where VGoT and VIC showed no significant correlation due to weak probe signals.
This validates the VLM-based binary detection as producing more reliable metrics than continuous CLIP probes.


%% ============================================================
\section{Bootstrap Confidence Intervals}
\label{app:bootstrap}
%% ============================================================

95\% bootstrap CIs ($B = 10{,}000$, percentile method) are reported in Table~\ref{tab:bootstrap_ci}.

\begin{table}[ht]
\caption{Main results with 95\% CIs.}
\label{tab:bootstrap_ci}
\centering
\small
\setlength{\tabcolsep}{2.5pt}
\begin{tabular}{@{}lcccc@{}}
\toprule
Metric & StoryDiff & EchoShot & VGoT & VIC \\
\midrule
PTM & \textbf{.243}{\tiny [.226,.262]} & .167{\tiny [.151,.184]} & .094{\tiny [.083,.107]} & .070{\tiny [.060,.082]} \\
EPA & \textbf{.593}{\tiny [.583,.604]} & .563{\tiny [.553,.574]} & .554{\tiny [.545,.564]} & .546{\tiny [.536,.557]} \\
\bottomrule
\end{tabular}
\end{table}

All PTM CIs are non-overlapping: SD [.226,.262] vs.\ Echo [.151,.184] vs.\ VGoT [.083,.107] vs.\ VIC [.060,.082].
EPA CIs are tight but also non-overlapping for all adjacent pairs, confirming statistically significant ordering.


%% ============================================================
\section{Statistical Significance and Effect Sizes}
\label{app:significance}
%% ============================================================

\begin{table}[ht]
\caption{Pairwise Mann-Whitney $U$ tests for PTM. Cohen's $d$ reports effect size.}
\label{tab:mann_whitney}
\centering
\small
\setlength{\tabcolsep}{2.5pt}
\begin{tabular}{@{}lccc@{}}
\toprule
Pair & $p$ & Cohen's $d$ & Sig. \\
\midrule
SD--Echo    & $3.3 \times 10^{-10}$ & 0.28 & $^{***}$ \\
Echo--VGoT  & $2.3 \times 10^{-11}$ & 0.32 & $^{***}$ \\
VGoT--VIC   & $2.9 \times 10^{-3}$  & 0.13 & $^{**}$  \\
SD--VIC     & $2.7 \times 10^{-53}$ & 0.72 & $^{***}$ \\
\bottomrule
\end{tabular}
\end{table}

Table~\ref{tab:mann_whitney} reports pairwise significance tests. All pairwise PTM differences are significant ($p < 0.01$).
The SD--VIC comparison has the largest effect size ($d = 0.72$), representing a 3.5$\times$ difference in PTM (0.243 vs.\ 0.070).
The VGoT--VIC comparison, while significant, has a smaller effect size ($d = 0.13$), consistent with both models being in the low-PTM regime.


%% ============================================================
\section{PTM Score Distribution}
\label{app:distributions}
%% ============================================================

\begin{table}[ht]
\caption{PTM score distributions. PTM takes discrete values $\{0, 0.5, 1.0\}$ for 2-transition patterns.}
\label{tab:ptm_dist}
\centering
\small
\begin{tabular}{@{}l cccc@{}}
\toprule
Model & $\mu \pm \sigma$ & $= 0$\% & $= 0.5$\% & $= 1.0$\% \\
\midrule
SD    & .243$\pm$.291 & 53.5 & 21.5 & 4.5 \\
Echo  & .167$\pm$.260 & 67.1 & 15.3 & 2.6 \\
VGoT  & .094$\pm$.197 & 79.5 & 10.2 & 0.8 \\
VIC   & .070$\pm$.173 & 84.8 & 8.7  & 0.3 \\
\bottomrule
\end{tabular}
\end{table}

Table~\ref{tab:ptm_dist} shows that PTM concentrates at discrete values because most patterns have 2 prescribed transitions, yielding PTM $\in \{0, 0.5, 1.0\}$.
VIC has 84.8\% of videos at PTM $= 0$ (no correct transitions) and only 0.3\% at PTM $= 1.0$ (all correct).
StoryDiffusion achieves PTM $= 1.0$ for 4.5\% of videos---rare but achievable.


%% ============================================================
\section{Frobenius Distance Analysis}
\label{app:frobenius}
%% ============================================================

The raw Frobenius distance $\|E\|_F$ is weakly discriminative (Table~\ref{tab:frobenius}):

\begin{table}[ht]
\caption{Pairwise Mann-Whitney $U$ tests for $\|E\|_F$.}
\label{tab:frobenius}
\centering
\small
\begin{tabular}{@{}lcc@{}}
\toprule
Pair & $p$-value & Mean$_1$ / Mean$_2$ \\
\midrule
SD--Echo    & 0.436 & 1.640 / 1.670 \\
SD--VGoT    & 0.051 & 1.640 / 1.613 \\
SD--VIC     & 0.363 & 1.640 / 1.643 \\
Echo--VGoT  & \textbf{0.004} & 1.670 / 1.613 \\
Echo--VIC   & 0.062 & 1.670 / 1.643 \\
VGoT--VIC   & 0.315 & 1.613 / 1.643 \\
\bottomrule
\end{tabular}
\end{table}

Only 1 of 6 pairs (Echo--VGoT) reaches significance at $\alpha = 0.05$, and the ranking ($\text{VGoT} < \text{SD} < \text{VIC} < \text{Echo}$) does \emph{not} match the PTM ranking.
This confirms that hold-cell noise dominates $\|E\|_F$, motivating the prescribed-subspace projection (PTM).


%% ============================================================
\section{Per-Pattern Failure Analysis}
\label{app:failure_pattern}
%% ============================================================

Failure rates are computed from the VLM-based $E$ matrix, with Type~IV (reversal, $|E|=2$) separated from Types~I/II (Table~\ref{tab:failure_rates}).

\begin{table}[ht]
\caption{Failure rates (\% of videos with $\geq$1 failure) per pattern, aggregated across all 4 models.}
\label{tab:failure_rates}
\centering
\small
\setlength{\tabcolsep}{3pt}
\begin{tabular}{@{}lccccc@{}}
\toprule
Pattern & $|\Delta S^* \neq 0|$ & I & II & III & IV \\
\midrule
Relay         & 2 & 63.4 & 68.6 & 34.9 & 11.5 \\
Split         & 3 & 73.1 & 70.8 & 37.3 & 13.4 \\
Accumulation  & 2 & 63.9 & 80.8 & 33.9 & 8.4 \\
Convergence   & 4 & 78.8 & 81.0 & 47.5 & 14.3 \\
Sliding Win.  & 4 & 77.3 & 77.5 & 37.3 & 12.5 \\
Reduction     & 2 & 63.3 & 68.8 & 28.5 & 8.0 \\
Rev.\ Relay   & 2 & 68.3 & 70.3 & 33.3 & 11.0 \\
\bottomrule
\end{tabular}
\end{table}

Patterns with more prescribed transitions (Convergence, Sliding Window: 4 each) show the highest failure rates across all types.
Amnesia (Type~II) dominates across all patterns, consistent with the entry--exit asymmetry (\S\ref{app:entry_exit}).


%% ============================================================
\section{GT Failure Validation}
\label{app:gt_failures}
%% ============================================================

To validate that our VLM detector and STEM pipeline correctly diagnose each failure type, we construct controlled GT videos from real cat footage.
We use two source clips---a black-and-white cat (entity~A) and an orange cat (entity~B)---plus a joint clip (A+B) to assemble five 3-shot videos following the Relay pattern ($S^* = [[1,0],[1,1],[0,1]]$, $\Delta S^* = [[0,{+}1],[{-}1,0]]$).

\paragraph{Setup.}
Each case substitutes specific source shots to produce a known $S_\text{obs}$:
\begin{itemize}[nosep,leftmargin=*]
\item \textbf{GT Perfect}: A $\to$ AB $\to$ B \quad $\Rightarrow$ PTM=1.000, EPA=1.000, $E=\mathbf{0}$
\item \textbf{Type~I (Context Bleeding)}: A $\to$ AB $\to$ AB \quad $\Rightarrow$ PTM=0.500, EPA=0.833
\item \textbf{Type~II (Amnesia)}: A $\to$ A $\to$ A \quad $\Rightarrow$ PTM=0.000, EPA=0.500
\item \textbf{Type~III (Spurious)}: B $\to$ AB $\to$ B \quad $\Rightarrow$ PTM=0.500, EPA=0.667
\item \textbf{Type~IV (Reversal)}: AB $\to$ A $\to$ B \quad $\Rightarrow$ PTM=0.500, EPA=0.667
\end{itemize}

\paragraph{Results.}
Running Qwen2-VL-7B-Instruct on the middle frame of each shot, all five cases produce $S_\text{obs}$ matrices that \emph{exactly} match the expected values.
The VLM achieves 100\% detection accuracy on these unambiguous, real-world images, confirming that:
(1)~the binary query protocol reliably detects clearly visible entities,
(2)~STEM correctly assigns the right failure type label in all cases, and
(3)~the lower TPR/TNR observed on generated images (\S\ref{app:vlm_quality}) reflects entity ambiguity in generation outputs, not a limitation of the detection pipeline.

Source videos and analysis scripts: \texttt{gt/two\_cats/} and \texttt{eval\_gt\_failures.py}.


%% ============================================================
\section{Metric Discriminability Comparison}
\label{app:metric_discriminability}
%% ============================================================

We compare the discriminative power of three candidate metrics: PTM ($\ell_0$-norm on prescribed cells), EPA (entity presence accuracy), and raw Frobenius distance $\|E\|_F$ (Table~\ref{tab:discriminability}).

\begin{table}[ht]
\caption{Discriminability of candidate metrics. $H$: Kruskal--Wallis statistic across 4 models. Sig.\ pairs: number of pairwise Mann--Whitney $U$ tests significant at $\alpha = 0.005$ (out of 6).}
\label{tab:discriminability}
\centering
\small
\begin{tabular}{@{}lccc@{}}
\toprule
Metric & $H$ & $p$ & Sig.\ pairs \\
\midrule
PTM              & \textbf{302.7} & $< 10^{-65}$ & \textbf{6/6} \\
EPA              & 46.0           & $5.5 \times 10^{-10}$ & 4/6 \\
$\|E\|_F$        & 8.9            & 0.031 & 1/6 \\
\bottomrule
\end{tabular}
\end{table}

PTM achieves 33$\times$ the $H$-statistic of $\|E\|_F$ and resolves all 6 pairwise comparisons, confirming that projecting onto the prescribed subspace ($E \odot M$) removes hold-cell noise that masks model differences.
EPA is intermediate (4/6 pairs), reflecting its role as a macroscopic complement rather than a primary metric.
The Frobenius norm resolves only the Echo--VGoT pair, and its ranking ($\text{VGoT} < \text{SD} < \text{VIC} < \text{Echo}$) does not match the PTM ranking---confirming it is confounded by unprescribed transitions.


%% ============================================================
\section{Sample Size Sensitivity}
\label{app:sample_size}
%% ============================================================

A key reviewer concern is whether 1,000 scenarios per model are necessary.
We subsample $N \in \{50, 100, 200, 300, 500, 750, 1000\}$ scenarios (with replacement, $B = 1{,}000$ bootstrap iterations) and measure two stability criteria (Table~\ref{tab:rank_stability}): (1)~\emph{rank preservation}: the fraction of bootstrap samples in which the 4-model PTM ranking matches the full-dataset ranking; (2)~\emph{full significance}: the fraction in which all 6 pairwise Mann--Whitney $U$ tests are significant at $p < 0.05$.

\begin{table}[ht]
\caption{Rank stability as a function of dataset size.}
\label{tab:rank_stability}
\centering
\small
\begin{tabular}{@{}ccc@{}}
\toprule
$N$ & Rank preserved (\%) & All 6 sig.\ (\%) \\
\midrule
50   & 72.8 & 21.3 \\
100  & 82.4 & 44.6 \\
200  & 89.7 & 72.1 \\
300  & 94.4 & 86.3 \\
500  & 97.6 & 96.5 \\
750  & 99.2 & 99.4 \\
1000 & 99.8 & 99.9 \\
\bottomrule
\end{tabular}
\end{table}

At $N = 300$, rank preservation exceeds 94\% and full significance exceeds 86\%, establishing a practical minimum for future benchmarks.
The 1,000-scenario design provides $>$99.8\% stability across both criteria, ensuring the reported results are robust to sampling variability.


%% ============================================================
\section{Cross-Pattern Rank Consistency}
\label{app:rank_consistency}
%% ============================================================

To verify that model rankings are not driven by a single easy/hard pattern, we compute per-pattern model rankings (Table~\ref{tab:rank_pattern}) and measure concordance via Kendall's coefficient of concordance $W$.

\begin{table}[ht]
\caption{Per-pattern PTM rankings (1 = best). Bold: deviation from overall ranking.}
\label{tab:rank_pattern}
\centering
\small
\setlength{\tabcolsep}{3pt}
\begin{tabular}{@{}lcccc@{}}
\toprule
Pattern & SD & Echo & VGoT & VIC \\
\midrule
Relay          & 1 & 2 & 3 & 4 \\
Split          & 1 & 2 & 3 & 4 \\
Accumulation   & 1 & 2 & 4 & 3 \\
Convergence    & 1 & 2 & 4 & 3 \\
Sliding Window & 1 & 2 & 3 & 4 \\
Reduction      & 1 & 2 & 3 & 4 \\
Reverse Relay  & 2 & \textbf{1} & 3 & 4 \\
\bottomrule
\end{tabular}
\end{table}

Kendall's $W = 0.869$ ($\chi^2 = 18.3$, $p = 3.9 \times 10^{-4}$), indicating strong concordance.
StoryDiffusion ranks 1st in 6/7 patterns; EchoShot claims 1st only on Reverse Relay (the simplest 2-transition derived pattern).
The VGoT/VIC swap in Accumulation and Convergence is consistent with VGoT's specific DynamiCrafter bottleneck on entry-heavy patterns (Appendix~\ref{app:entry_exit}).


%% ============================================================
\section{Reproducibility Checklist}
\label{app:reproducibility}
%% ============================================================

\begin{itemize}[leftmargin=*,nosep]
\item All evaluation code, prompts, and 4,000 generated videos are released.
\item Fixed random seed for generation and ablation bootstraps.
\item VLM detector: Qwen2-VL-7B-Instruct (fp16, greedy decoding).
\item Entity query: ``Is there a \{entity\} visible in this image? Answer only Yes or No.''
\item 1 middle frame per shot (3 frames per video).
\item PTM: prescribed transitions only ($\Delta S^* \neq 0$), no threshold.
\item EPA: exact match $S_\text{obs}[t,k] = S^*[t,k]$, no threshold.
\item Bootstrap: $B = 10{,}000$ iterations, percentile CIs.
\item Permutation test: 1,000 character role swappings.
\item Hardware: 4$\times$ NVIDIA RTX 4090 (24 GB each).
\item Evaluation runtime: $\sim$20 minutes for 4,000 videos (VLM inference).
\item Table reproduction: \texttt{generate\_tables.py} reproduces all paper tables from \texttt{vlm\_fullscale\_merged.json}.
\end{itemize}


\bibliographystyle{ACM-Reference-Format}
\bibliography{references}

\end{document}